\section{Baseline Details and References}
\label{ap:baselines}

Due to the broad range of baselines considered in our experiments, we provide additional details in this appendix, including how each baseline is adapted to our evaluation settings when necessary. For complete methodological descriptions, we refer readers to the corresponding original publications.


\paragraph{Traditional Time Series Models.}
This group includes classical statistical models and neural models that operate directly on numerical time-series histories. These methods serve as non-agentic baselines for evaluating how far standard temporal modeling can go without explicit contextual reasoning or agentic workflow support.

\begin{itemize}
   \item \textbf{ARIMA}.
    ARIMA is a classical statistical model that captures temporal dependence through autoregressive, differencing, and moving-average components.


    \item \textbf{ETS}.
    ETS uses exponential smoothing to model level, trend, and seasonal
    components. Like ARIMA, it relies on historical numerical patterns rather than natural-language context or agentic reasoning.

    \item \textbf{DLinear}~\citep{DBLP:conf/aaai/ZengCZ023}.
    DLinear decomposes a time series into trend and remainder components using a moving average, applies separate one-layer linear mappings to them, and sums the outputs for prediction.

    \item \textbf{PatchTST}~\citep{DBLP:conf/iclr/NieNSK23}.
    PatchTST is a Transformer-based model for long-term time-series modeling. It segments each time series into subseries-level patches and treats these patches as input tokens to the Transformer. It also uses a channel-independent design, where different variables share the same embedding and Transformer weights.
\end{itemize}

\paragraph{Time Series Foundation Models.}
This group contains pretrained foundation models designed for broad
time-series modeling across datasets and domains. These baselines test whether specialized temporal pretraining alone can match a harnessed agent that uses temporal tools, contextual reasoning, reusable routines, and memory.

\begin{itemize}
    \item \textbf{Chronos-family models}~\citep{DBLP:journals/tmlr/AnsariSTZMSSRPK24, DBLP:journals/corr/abs-2510-15821}.
    Chronos formulates time-series forecasting as language modeling over discrete value tokens: continuous observations are scaled, quantized into a fixed vocabulary, and modeled with Transformer language-model architectures. Through pretraining on diverse real and synthetic time series, Chronos-family models provide strong zero-shot probabilistic forecasting baselines.

    \item \textbf{Lag-Llama}~\citep{rasul2023lag}.
    Lag-Llama is a pretrained foundation model for univariate probabilistic time-series forecasting. It uses a decoder-only Transformer architecture with lagged values as covariates and is pretrained on diverse time-series datasets to support zero-shot and fine-tuned prediction. 
    \item \textbf{Moirai-family models}~\citep{DBLP:conf/icml/WooLKXSS24}.
    Moirai is a universal time-series forecasting Transformer trained for cross-domain zero-shot prediction. It is designed to handle heterogeneous temporal data, including different sampling frequencies, variable counts, and distributional properties, through a masked encoder architecture.
\end{itemize}

\paragraph{Language Models.}
This group includes general-purpose language models and time-series-oriented language-model baselines. These methods help evaluate whether language modeling or language-based adaptation alone is sufficient for contextualized temporal reasoning, without the full time-series-native harness used by \method. In our experiments, we compared with a broader set of such general-purpose language models show improved performance over them. We provide details of representative models below. We refer the readers to the original benchmark papers \cite{DBLP:conf/icml/WilliamsAMZSRRL25, DBLP:journals/corr/abs-2601-18744} for more details.

\begin{itemize}
    \item \textbf{LLaMA3-70B}~\citep{DBLP:journals/corr/abs-2407-21783}.
    LLaMA3-70B is a strong open-weight general-purpose LLM from the Llama 3 family. For CiK, we use it as a text-only language-model baseline: the numerical time series and natural-language context are serialized into a single prompt, and the model generates the continuation autoregressively through next-token prediction. Under greedy decoding, each generated token is selected according to the model's output logits.

    \item \textbf{UniTime}~\citep{DBLP:conf/www/LiuHLDLHZ24}.
    UniTime is a language-empowered unified model for cross-domain time-series forecasting. It uses domain instructions to provide domain-level information and a Language-TS Transformer to align language and time-series representations. The model aims to handle heterogeneous time series from different domains under
    a unified architecture.

    \item \textbf{Time-LLM}~\citep{DBLP:conf/iclr/0005WMCZSCLLPW24}.
    Time-LLM reprograms pretrained LLMs for time-series forecasting while keeping the backbone language model frozen. It maps time-series patches into text-prototype representations aligned with the LLM input space, augments the input with Prompt-as-Prefix task information, and projects the LLM outputs into temporal predictions.
\end{itemize}

\paragraph{Agentic Frameworks.}
This group includes general prompting and agentic workflows that are not specifically designed for time series. They test whether generic LLM reasoning, tool-use, or reflection pipelines are sufficient for contextualized temporal tasks without specialized temporal runtime support.

\begin{itemize}
    \item \textbf{Direct prompting.}
    Direct prompting feeds the serialized time series, context, and task instruction to the language model in a single prompt. The model directly produces the final answer without explicit reasoning steps, tool use, reflection, or multi-agent interaction.

    \item \textbf{Chain-of-thought prompting}~\citep{DBLP:conf/nips/Wei0SBIXCLZ22}.
    Chain-of-thought prompting asks the language model to produce intermediate natural-language reasoning steps before the final answer. It is a prompting-only strategy for eliciting multi-step reasoning from large language models.

    \item \textbf{ReAct}~\citep{DBLP:conf/iclr/YaoZYDSN023}. 
    ReAct interleaves natural-language reasoning traces with task-specific actions. The reasoning steps help the model maintain and update its plan, while actions allow it to interact with external tools or environments. It is a general agentic prompting baseline for combining reasoning and tool use.

    \item \textbf{Self-reflection}~\citep{DBLP:journals/corr/abs-2405-06682}.
    Self-reflection lets an LLM inspect its previous answer or reasoning process, identify possible errors, and use the generated feedback to revise its response. The authors study several reflection variants where agents learn guidance from prior mistakes before attempting the task again.

    \item \textbf{Multi-agent reflection}~\citep{DBLP:conf/icml/YuanX25}.
    Multi-agent reflection uses multiple LLM agents to iteratively verify and
improve answers. It models the refinement process as a Markov Decision Process and trains an actor-critic LLM system to incorporate feedback over multiple rounds, improving reasoning through coordinated agent collaboration.
\end{itemize}

\paragraph{Time Series Agents.}
This group includes agentic methods specifically designed for time-series
analysis. These baselines are closest to our setting, but they typically focus
on particular temporal tasks or fixed agentic workflows rather than providing a
general harness for contextualized temporal reasoning.

\begin{itemize}
    \item \textbf{TS-Agent}~\citep{DBLP:journals/corr/abs-2510-07432}.
    TS-Agent is a time-series reasoning agent that lets LLMs synthesize conclusions from evidence while delegating statistical and structural extraction to tools over raw numerical sequences. It maintains an explicit evidence log and uses a self-critic with a final quality gate to iteratively refine reasoning.

    \item \textbf{TSci}~\citep{DBLP:journals/corr/abs-2510-01538}.
    TSci is an LLM-driven agentic framework for general time-series forecasting. It organizes the workflow into specialized agents for data diagnostics, preprocessing, model planning, fitting, validation, ensembling, and reporting. The framework emphasizes automated forecasting together with transparent natural-language rationales and reports.
\end{itemize}